\documentclass[11pt]{article}
\usepackage[T1]{fontenc}
\usepackage{lmodern}
\usepackage[margin=1in]{geometry}
\usepackage{amsmath,amssymb,amsthm,mathtools}
\usepackage{microtype}
\usepackage{booktabs}
\usepackage{xcolor}
\usepackage{enumitem}
\usepackage{hyperref}
\hypersetup{colorlinks=true,linkcolor=blue!45!black,citecolor=blue!45!black,
 urlcolor=blue!45!black,filecolor=blue!45!black,
 pdftitle={Arithmetic obstructions in positive determinant approximations to beta(4)},
 pdfauthor={Research note prepared with Codex}}
\urlstyle{same}
\setlist{nosep}
\setlength{\emergencystretch}{1em}
\newtheorem{theorem}{Theorem}[section]
\newtheorem{proposition}[theorem]{Proposition}
\newtheorem{lemma}[theorem]{Lemma}
\newtheorem{corollary}[theorem]{Corollary}
\theoremstyle{definition}
\newtheorem{definition}[theorem]{Definition}
\theoremstyle{remark}
\newtheorem{remark}[theorem]{Remark}
\numberwithin{equation}{section}
\newcommand{\Q}{\mathbb Q}
\newcommand{\Z}{\mathbb Z}
\newcommand{\R}{\mathbb R}
\newcommand{\N}{\mathbb N}
\newcommand{\dd}{\,\mathrm d}
\newcommand{\vp}{v_p}
\newcommand{\coeff}[2]{[X^{#1}]#2}
\newcommand{\PP}{\operatorname{Prim}}
\DeclareMathOperator{\sech}{sech}
\DeclareMathOperator{\diag}{diag}
\DeclareMathOperator{\Ree}{Re}
\title{Arithmetic obstructions in positive determinant\newline
approximations to $\beta(4)$}
\author{\normalsize Research note prepared with Codex}
\date{31 October 2026}

\begin{document}
\maketitle

\begin{abstract}
Small positive determinants offer a possible route to proving that
$\beta(4)=1-3^{-4}+5^{-4}-\cdots$ is irrational. The determinant must
be the value of an integer polynomial, and it must be sufficiently
small relative to the polynomial's degree. We examine two explicit
constructions and prove that clearing denominators necessarily destroys
the required smallness, even after all common coefficient factors are
removed. The resulting primitive integer-polynomial values
grow at least exponentially in the square of the matrix size, with
lower exponents $2-2\log2$ and $6-\log4$. The first conclusion covers
every consecutive block of positive integer nodes, regardless of how
its starting point varies with the size. For the second construction,
we determine the exact contribution of every sufficiently large prime
to the primitive multiplier. These are failure theorems for the two
specified families. We explain the irrationality criterion and the
constructions before developing their positive integral representations,
real estimates and prime-divisibility arguments.
\end{abstract}

\section{Introduction}

We study two attempts to prove that
\begin{equation}\label{eq:theta}
 \theta=\beta(4)=\sum_{k=0}^{\infty}\frac{(-1)^k}{(2k+1)^4}
\end{equation}
is irrational. Both attempts produce positive determinants that become
very small. The question is whether these small numbers can be turned
into a contradiction to rationality. We prove that, for these two
constructions, the conversion to integer coefficients necessarily
destroys the smallness. Theorems~\ref{thm:principal} and~\ref{thm:delta}
are thus results about the limitations of particular proof strategies.

\subsection{What would prove irrationality?}

Suppose that $\theta=a/b$ were rational, with integers $a,b$ and $b>0$.
For any integers $A,B$,
\[
 B\theta-A=\frac{Ba-Ab}{b}.
\]
If this is nonzero, its numerator is a nonzero integer, so its absolute
value is at least $1/b$. Consequently, a sequence of nonzero integer
linear forms $B_n\theta-A_n$ tending to zero would prove irrationality:
no fixed denominator $b$ could accommodate the whole sequence.

This is stronger than finding accurate rational approximations.
If $A_n/B_n$ approximates $\theta$, the relevant quantity is
\[
 |B_n\theta-A_n|=|B_n|\left|\theta-\frac{A_n}{B_n}\right|.
\]
The approximation error must remain small after multiplication by its
denominator. Every real number has arbitrarily accurate rational
approximations; the integer-form requirement is what makes this an
irrationality test.

\subsection{Why use determinants?}

If the entries in a matrix are rational linear forms in $\theta$, its
determinant is a polynomial with rational coefficients evaluated at
$\theta$. For example, with $r,s,t\in\Q$,
\[
 \det\begin{pmatrix}X+r&s\\s&X+t\end{pmatrix}
       =X^2+(r+t)X+rt-s^2.
\]
An $n\times n$ matrix of linear entries similarly gives a polynomial
of degree at most $n$, and exactly $n$ in our two constructions.
Increasing the dimension can help make the determinant small, but
increasing the degree makes the irrationality test harder.

On the \emph{analytic} side, we estimate the determinant's absolute
value and how rapidly it decreases as the matrix size grows. This
part of the argument treats the determinant as a real number and does
not assume that $\theta$ is rational. Determinants can be especially
small when their rows are almost linearly dependent.
Orthogonalizing the rows of a real
square matrix by subtracting combinations of earlier rows preserves
the determinant. Its absolute value is the product of the lengths of
the orthogonalized row vectors. Thus small components in successive new directions can
multiply together. For illustration, lengths
$e^{-c},e^{-2c},\ldots,e^{-nc}$, with $c>0$, would give
\[
 |\det M_n|=\prod_{j=1}^n e^{-cj}
          =e^{-cn(n+1)/2}.
\]
This explains how decay on the $n^2$ scale can arise. It is an
illustrative mechanism; the rate for each actual construction requires
proof, as does nonvanishing.

On the \emph{arithmetic} side, we use properties of integers and
fractions. The starting fact is that a nonzero integer has absolute
value at least one. This gives a lower bound for a nonzero rational
number once we control its denominator. Suppose $\theta=a/b$ and let
$Z\in\Z[X]$, meaning that $Z$ is a polynomial with integer coefficients.
If its degree is at most $n$, multiplying $Z(a/b)$ by $b^n$ clears
all denominators. Hence $b^nZ(a/b)$ is an integer, and
\begin{equation}\label{eq:integer-test}
 Z(\theta)\ne0\quad\Longrightarrow\quad
 |Z(\theta)|\geq b^{-n}\qquad\text{if }\theta=a/b.
\end{equation}
Larger degree permits smaller nonzero values at rational arguments;
this is a cost of the construction. The potential gain is that the
determinant can shrink on the $n^2$ scale while the degree grows only
linearly. If, after conversion to integer coefficients, its nonzero
value is at most $e^{-cn^2}$ for some fixed $c>0$, then it is eventually
smaller than $b^{-n}$ for every fixed $b$, contradicting rationality.
More generally, it suffices that $|Z_n(\theta)|^{1/n}\to0$.
Merely tending to zero is insufficient:
$Z_n(X)=X^n$ has value $2^{-n}$ at the rational number $1/2$.

Positive determinant constructions are attractive because they can
meet two requirements at once: rapid decay and nonvanishing. A
\emph{Gram matrix} consists of the pairwise inner products of a list
of functions. If those functions are linearly independent, the matrix
is positive definite and its determinant is positive, although it may
be very small. Integral representations of this kind let us establish
nonvanishing without knowing whether $\theta$ is rational.
Zudilin~\cite{Zudilin} develops this approach systematically.
The determinant strategy for the related value $\beta(2)$, Catalan's
constant, described by OpenAI~\cite[Section~1.2]{OpenAI}, is a motivation
for testing such constructions at $\beta(4)$.

\subsection{The cost of integer coefficients}

The determinant initially gives $Q_n\in\Q[X]$, whereas the test
\eqref{eq:integer-test} requires integer coefficients. We therefore
clear denominators and then divide out the greatest common divisor of
all coefficients. The resulting integer polynomial is called
\emph{primitive}. Write
\begin{equation}\label{eq:budget}
 Z_n=\kappa_n Q_n=\PP(Q_n),\qquad
 \log|Z_n(\theta)|=\log\kappa_n+\log|Q_n(\theta)|,
 \quad \kappa_n>0.
\end{equation}
Every integer polynomial proportional to $Q_n$ is an integer multiple
of $Z_n$. Thus $Z_n$ has the smallest nonzero absolute value obtainable
by rationally rescaling this same polynomial.

Removing common factors matters. For example,
$Q(X)=(6X+12)/5$ becomes integral after multiplication by $5$, but its
primitive representative is only $X+2$: the actual multiplier is $5/6$.
A large denominator chosen to clear the entries of a matrix could
therefore be a poor estimate of the cost of its determinant.
Brown~\cite[Section~7.4]{Brown} discusses the role of common divisors
and congruences in this arithmetic comparison.

To prove that a construction fails, we must bound the actual primitive
value from below. An upper bound that fails to tend to zero establishes
only that the estimate is inadequate. Our proofs instead show that
the determinants cannot be smaller, or their primitive multipliers
cheaper, by enough to make the construction work.

\subsection{The two results and what they mean}

The exact matrix entries are given in
Sections~\ref{sec:kernel-construction} and~\ref{sec:gram-construction}.
Here we describe the two designs and state the conclusions.
Carreira~\cite{Carreira} also investigates $\beta(4)$ determinant
constructions and their primitive arithmetic. The results below give
proofs of failure for the specified infinite families, including all
common coefficient cancellation.

\paragraph{First construction: consecutive nodes.}
We build a symmetric matrix with rational off-diagonal entries and
$\theta$ plus a rational on each diagonal entry. The same nodes are
used for its rows and columns; this is the meaning of \emph{principal}.
We choose consecutive integer nodes
\[
 m+1,m+2,\ldots,m+n.
\]
The parameter $n$ is the matrix size and $m$ is its offset. Replacing
$\theta$ by $X$ and taking the determinant defines $Q_{n,m}(X)$;
write $Z_{n,m}=\PP(Q_{n,m})$. The construction is made explicit in
\eqref{eq:R}--\eqref{eq:principal-Q}.

\begin{theorem}[Consecutive principal determinants]\label{thm:principal}
For every sequence of nonnegative integer offsets $m_n$,
$Q_{n,m_n}(\theta)>0$ and
\begin{equation}\label{eq:principal-result}
 \liminf_{n\to\infty}\frac{\log Z_{n,m_n}(\theta)}{n^2}
 \geq 2-2\log 2>0.
\end{equation}
\end{theorem}

The positive constant $2-2\log2$ is approximately $0.614$.
The inequality says that, for any chosen offset sequence $m_n$ and
any $\varepsilon>0$, eventually
\[
 Z_{n,m_n}(\theta)\geq
       e^{(2-2\log2-\varepsilon)n^2}.
\]
These integer-polynomial values grow instead of becoming small.
The offset may depend on $n$ in any way, even growing much faster
than $n$. Thus moving the consecutive block farther out cannot repair
this family. The theorem concerns consecutive blocks, not arbitrary
selections of nodes.

\paragraph{Second construction: moments of a positive weight.}
We use integrals of powers of $y=(u-v)^2$, with $u,v\in[0,1]$, against
a positive weight. Such integrals are called \emph{moments}. Write
them as $\mu_0,\mu_1,\ldots$ and form the matrix
$[\mu_{i+j}]_{0\leq i,j<n}$, the Gram matrix of
$1,y,\ldots,y^{n-1}$. Each moment turns out to be a rational multiple
of $\theta$ plus a rational. Replacing $\theta$ by $X$ gives a
determinant polynomial $Q_n^\Delta(X)$ and its primitive representative
$Z_n^\Delta=\PP(Q_n^\Delta)$. Formula~\eqref{eq:delta-Q} gives their
precise definitions; the superscript $\Delta$ denotes this
difference-variable construction.

\begin{theorem}[Difference-variable Gram determinants]\label{thm:delta}
The polynomial $Q_n^{\Delta}$ has degree $n$ and leading coefficient
$2^{n-1}$, and $Q_n^{\Delta}(\theta)>0$. Moreover,
\begin{equation}\label{eq:delta-result}
 \log Z_n^{\Delta}(\theta)\geq (6-\log 4)n^2+o(n^2).
\end{equation}
\end{theorem}

Here $o(n^2)$ denotes a term whose ratio to $n^2$ tends to zero.
Since $6-\log4\approx4.614>0$, these primitive values also grow
exponentially in $n^2$. In this family we can say more about the
reason: Proposition~\ref{prop:gram-prime-cost} determines exactly
how many powers of every prime $p>4n$ remain in the primitive
multiplier. That arithmetic result is stated alongside its proof
in Section~\ref{sec:delta-arithmetic}.

These conclusions go beyond observing that a list of examples fails.
They rule out the required small values at all sufficiently large
sizes, with every offset schedule included in the first family.
They also account for all rational scalar rescalings, because the
primitive polynomial already gives the smallest such value.
The contribution is a precise explanation of failure for these two
natural constructions. It does not decide whether $\theta$ is rational,
or exclude different determinants, polynomial factors or combinations
of the constructed expressions.

\subsection{How the proofs work}

The two proofs compare the gain from a small real determinant with
the cost of making its coefficients primitive integers. Both use
the same three steps:
\begin{enumerate}
\item Express the determinant through a positive integral, so its
      value cannot be zero.
\item Bound that value both above and below, determining its rate
      of decay. A lower bound is essential for a failure theorem.
\item Find prime powers that must occur in the primitive multiplier,
      even after every common coefficient factor has been removed.
\end{enumerate}

The comparison is particularly simple at the leading quadratic scale.
In the table, each entry is the coefficient of $n^2$ in the indicated
logarithm; terms $o(n^2)$ are omitted. The first row uses offset $m=0$.
\begin{center}
\begin{tabular}{lccc}
\toprule
Family & $\log Q_n(\theta)$ & $\log\kappa_n$ (lower)
       & $\log Z_n(\theta)$ (lower)\\
\midrule
Consecutive, $m=0$ & $-2\log2$ & $2$ & $2-2\log2$\\
Difference Gram & $-\log4$ & $6$ & $6-\log4$\\
\bottomrule
\end{tabular}
\end{center}
The negative second column describes the small determinants. Adding
the arithmetic cost in the third column gives a positive fourth
column: the best integer-coefficient values become large.
The proof of Theorem~\ref{thm:principal} also checks that changing
the offset cannot reverse this comparison.

Section~\ref{sec:normalization} gives the arithmetic conventions.
Sections~\ref{sec:kernel}--\ref{sec:principal} construct and analyze
the first family; Sections~\ref{sec:gram}--\ref{sec:delta-arithmetic}
do so for the second. Each construction is introduced before its
technical proof begins, and the two arguments can be read separately.
Section~\ref{sec:scope} describes the scope and evidence; Appendix~\ref{app:zero}
includes a zero node in the first construction.

\begin{samepage}
The introductory argument uses rational arithmetic and linear algebra.
The detailed proofs also use calculus, elementary complex integration,
Stirling's formula and the prime number theorem in the form
\begin{equation}\label{eq:PNT}
 \vartheta(x):=\sum_{p\leq x}\log p=x+o(x).
\end{equation}
\end{samepage}
The latter counts the total logarithmic contribution of the primes
whose powers the arithmetic arguments force. No numerical optimization
or finite enumeration is an input to either asymptotic obstruction.

\section{Primitive normalization and prime powers}\label{sec:normalization}

\begin{definition}
For $0\ne Q\in\Q[X]$, its primitive representative $\PP(Q)$ is the
unique positive rational multiple $\kappa Q\in\Z[X]$ whose coefficients
have gcd one. This convention preserves the sign of $Q$.
\end{definition}

We write $[X^j]Q$ for the coefficient of $X^j$ in $Q$.
For a nonzero rational number, $\vp$ counts powers of $p$ in its
numerator minus those in its denominator; set $\vp(0)=+\infty$.
For example $v_p(p^{-3})=-3$. Saying that a rational
number is $p$-integral means that its denominator in lowest terms is not
divisible by $p$. Reduction modulo $p$ below is always performed on
$p$-integral quantities, by inverting the denominator modulo $p$.
A $p$-unit has valuation zero, so its reduction is nonzero.

\begin{lemma}\label{lem:primitive}
If $Q(X)=\sum_j q_jX^j\ne0$ and $\PP(Q)=\kappa Q$, then
\begin{equation}\label{eq:min-valuation}
 \vp(\kappa)=-\min_j\vp(q_j)
\end{equation}
for every prime $p$. If $Q$ is monic, $\kappa$ is a positive integer.
If the leading coefficient of $Q$ is $2^e$ with $e\geq0$, then
$\vp(\kappa)\geq0$ for odd $p$ and $v_2(\kappa)\geq-e$.
\end{lemma}

\begin{proof}
Integral primitive coefficients have minimum valuation zero at every
prime, which proves~\eqref{eq:min-valuation}. A monic polynomial has a
coefficient of valuation zero at every prime. A leading coefficient $2^e$
has valuation zero at odd primes and valuation $e$ at 2.
\end{proof}

If $Z=\PP(Q)$ and $\lambda Q$ has integer coefficients, then
$(\lambda/\kappa)Z$ has integer coefficients. The denominator of
$\lambda/\kappa$ in lowest terms must divide every coefficient of $Z$,
so it is one. Thus every integral rational rescaling is an integer
multiple of $Z$, as asserted in the introduction.

Formula~\eqref{eq:min-valuation} tells us how to prove that a prime
power is unavoidable: to show $v_p(\kappa)\geq r$, it suffices to find
one coefficient of $Q$ with valuation at most $-r$. To establish equality,
we must also show that no coefficient has a lower valuation.
The first construction uses the lower bound; the second obtains
equality for every sufficiently large prime.

If $M(X)$ has rational polynomial entries, $\det M(X)\ne0$, and $U,V$
are invertible rational matrices of the same size, then
$\det(U M(X)V)=\det U\det V\det M(X)$. Consequently the primitive
representative of $\det M(X)$ is unchanged up to sign by rational invertible
row and column operations. Passing to a different subspace or a different
matrix is a different construction and is not covered by this observation.

\section{First construction: a positive kernel}\label{sec:kernel}

This section constructs the first family from a shifted beta-series
tail. We then prove positivity and estimate the real determinant.
Section~\ref{sec:principal} will determine the arithmetic cost.

\subsection{Choosing the matrix entries}\label{sec:kernel-construction}

The first construction seeks a symmetric matrix with rational
off-diagonal entries and $\theta$ plus a rational on each
diagonal entry. We also want the matrix to be positive definite, so
that its determinant cannot vanish. The following calculation explains
the entries; the integral representation in the next subsection will
establish positive definiteness.

For integers $a\geq 0$ and $\ell\geq 2$, write
\begin{equation}\label{eq:S}
 S_\ell(a)=\sum_{j=0}^{a-1}
 \left(\frac1{(4j+1)^\ell}-\frac1{(4j+3)^\ell}\right),
 \qquad S_\ell(0)=0.
\end{equation}
These are the first $2a$ terms of the alternating beta series of order
$\ell$, grouped into pairs. In particular, $\theta-S_4(a)$ is the
remaining fourth-power tail. To obtain it by differentiation, start
from the shifted cubic tail
\[
 f(t)=\sum_{k=0}^{\infty}\frac{(-1)^k}{(2k+1+4t)^3}
 \qquad(t\geq0).
\]
At a nonnegative integer $a$, shifting the summation index and
differentiating give
\[
 f(0)-f(a)=S_3(a),\qquad f'(a)=-12\bigl(\theta-S_4(a)\bigr).
\]
Thus subtracting the tail from $f(0)$ gives rational values at integer
arguments, while its derivative introduces $\theta$. Combine these
properties by setting, for $x>0$,
\[
 g(x)=\frac{\sqrt{x}}6\bigl(f(0)-f(\sqrt{x})\bigr),\qquad
 K(x,y)=\frac{g(y)-g(x)}{y-x},\quad K(x,x)=g'(x).
\]
Here $K$ is a divided difference, with its limiting derivative on the
diagonal. At squared integer nodes,
\[
 g(a^2)=\frac{aS_3(a)}6,\qquad
 g'(a^2)=\frac{S_3(a)-af'(a)}{12a}
        =\theta-S_4(a)+\frac{S_3(a)}{12a}\quad(a>0).
\]
The factor $\sqrt{x}$ compensates for the $1/(2\sqrt{x})$ from the
chain rule, and the factor $1/6$ makes the coefficient of $\theta$
exactly one. These identities explain both the squared nodes and
the constants in the matrix entries. For distinct positive integers
$a_1,\ldots,a_n$, we have
$[K(a_i^2,a_j^2)]_{i,j}=\theta I_n+R(a)$, where $R(a)$ is the rational
symmetric matrix
\begin{equation}\label{eq:R}
 R(a)_{ii}=-S_4(a_i)+\frac{S_3(a_i)}{12a_i},\qquad
 R(a)_{ij}=\frac{a_jS_3(a_j)-a_iS_3(a_i)}{6(a_j^2-a_i^2)}\quad(i\ne j).
\end{equation}
The integral representation below proves that $K$ is a
positive kernel: its matrices at distinct positive nodes are positive
definite. This does not follow from the divided-difference formula
alone. Replacing $\theta$ by $X$ gives a \emph{matrix pencil} $XI_n+R(a)$,
whose entries depend linearly on $X$ and whose determinant is monic.
For consecutive nodes $a_i=m+i$, $1\leq i\leq n$, set
\begin{equation}\label{eq:principal-Q}
 Q_{n,m}(X)=\det(XI_n+R(a)),\qquad
 Z_{n,m}(X)=H_{n,m}Q_{n,m}(X)=\PP(Q_{n,m})(X),
\end{equation}
where $H_{n,m}>0$ and $Z_{n,m}\in\Z[X]$ has coefficient gcd one.

To obtain positivity and the size estimate, we express $K(x,y)$ as
an integral of $(x+s^2)^{-1}(y+s^2)^{-1}$ against a positive density.
These factors give a Gram representation when the row and column nodes
agree. An explicit determinant integral handles different node lists.
We then compare the density with simpler ones that yield Cauchy
determinants.

\subsection{A direct integral representation}

Extend the preceding real-variable function to $\Ree z\geq0$ by
\begin{equation}\label{eq:f}
 f(z)=\sum_{j=0}^{\infty}
 \left((4j+1+4z)^{-3}-(4j+3+4z)^{-3}\right).
\end{equation}
Expanding $\sech u=2\sum_{k\geq0}(-1)^ke^{-(2k+1)u}$ and integrating
gives
\begin{equation}\label{eq:f-laplace}
 f(z)=\frac14\int_0^\infty u^2\sech(u)e^{-4zu}\dd u,
 \quad f(a)=f(0)-S_3(a),\quad f'(a)=-12(\theta-S_4(a))
\end{equation}
at nonnegative integer $a$. The integrated series is absolutely summable,
which justifies the expansion and differentiation used here.

The next two transforms convert the exponential in this Laplace
integral into the rational factors needed for the kernel. The cosine
transform evaluates the resulting density explicitly; the second
identity replaces $\sqrt{x}(1-e^{-t\sqrt{x}})$ by an integral
containing $x/(x+s^2)$. We include their proofs to fix the normalization:
\begin{align}
 \int_0^\infty\sech(u)\cos(tu)\dd u
   &=\frac\pi2\sech\frac{\pi t}{2},\label{eq:sech-transform}\\
 \int_0^\infty\frac{x(1-\cos(ts))}{x+s^2}\dd s
   &=\frac{\pi\sqrt{x}}2(1-e^{-t\sqrt{x}}),
       \qquad x>0,\ t\geq0.\label{eq:cos-transform}
\end{align}
For~\eqref{eq:sech-transform}, integrate $e^{itz}/\cosh z$ around the
rectangle with vertices $\pm R,\pm R+i\pi$. The top edge contributes
$e^{-\pi t}$ times the bottom edge, and the pole at $i\pi/2$ has residue
$e^{-\pi t/2}/i$. The vertical edges tend to zero. Hence the full-line
transform is $\pi\sech(\pi t/2)$; symmetry gives the displayed formula.
For~\eqref{eq:cos-transform}, the upper-half-plane pole at $i\sqrt{x}$
gives
$\int_0^\infty\cos(ts)/(s^2+x)\dd s
 =\pi e^{-t\sqrt{x}}/(2\sqrt{x})$ for $t>0$, and the value at $t=0$
follows by continuity. Subtract the latter transforms and multiply by $x$.

Differentiating~\eqref{eq:sech-transform} twice under the integral gives
\begin{equation}\label{eq:f-fourier}
 f(0)=\frac{\pi^3}{32},\qquad
 \Ree f(is)=\frac{\pi^3}{32}A(s)(2A(s)^2-1),
 \quad A(s)=\sech(2\pi s).
\end{equation}
Define
\begin{equation}\label{eq:gK}
 g(x)=\frac{\sqrt{x}}6\bigl(f(0)-f(\sqrt{x})\bigr),\quad g(0)=0,
 \qquad K(x,y)=\frac{g(y)-g(x)}{y-x},
\end{equation}
with the derivative value when $x=y$.

\begin{proposition}\label{prop:kernel}
For $x,y\geq0$,
\begin{equation}\label{eq:K-integral}
 K(x,y)=\int_0^\infty
 \frac{\omega(s)}{(x+s^2)(y+s^2)}\dd s,
 \quad
 \omega(s)=\frac{\pi^2}{96}s^2q(s),\quad
 q(s)=1+A(s)-2A(s)^3.
\end{equation}
The density $\omega$ is strictly positive for $s>0$. On squared integer
nodes the kernel equals
\begin{align}
 K(a^2,b^2)&=\frac{bS_3(b)-aS_3(a)}{6(b^2-a^2)}&& (a\ne b),\label{eq:K-off}\\
 K(a^2,a^2)&=\theta-S_4(a)+\frac{S_3(a)}{12a}&& (a>0),\label{eq:K-diag}\\
 K(0,0)&=2\theta.\label{eq:K-zero}
\end{align}
Every square minor with strictly increasing nonnegative row nodes and
strictly increasing nonnegative column nodes is positive.
\end{proposition}

\begin{proof}
First identify the density. Equations~\eqref{eq:f-laplace} and
\eqref{eq:cos-transform}, with $t=4u$, give
\[
 g(x)=\frac1{12\pi}\int_0^\infty\int_0^\infty
 u^2\sech u\,\frac{x(1-\cos(4us))}{x+s^2}\dd s\dd u.
\]
The integrand is nonnegative, so Tonelli's theorem permits interchange.
The inner $u$-integral is $4(f(0)-\Ree f(is))$. Hence
\begin{equation}\label{eq:g-stieltjes}
 g(x)=\int_0^\infty\frac{x}{x+s^2}\sigma(s)\dd s,
 \qquad \sigma(s)=\frac{f(0)-\Ree f(is)}{3\pi}
                =\frac{\pi^2}{96}q(s).
\end{equation}
Taking a divided difference uses the elementary identity
\[
 \frac{y/(y+s^2)-x/(x+s^2)}{y-x}
       =\frac{s^2}{(x+s^2)(y+s^2)},
\]
which gives~\eqref{eq:K-integral}. Positivity is now explicit:
\begin{equation}\label{eq:q-positive}
 q=(1-A)(2A^2+2A+1)>0\qquad(s>0).
\end{equation}
Since $\omega(s)=O(s^4)$ at zero and $O(s^2)$ at infinity, the integral
is finite even when $x=y=0$. Differentiation or dominated convergence
justifies the diagonal and zero limits. The identities
\eqref{eq:K-off}--\eqref{eq:K-diag} follow from~\eqref{eq:f-laplace};
at zero, $g(x)=-f'(0)x/6+O(x^{3/2})=2\theta x+O(x^{3/2})$.

It remains to pass from the kernel to its minors. For increasing
nodes $x_i,y_j$, expanding the two determinants and
integrating termwise gives the determinant integration identity
\begin{equation}\label{eq:andreief}
 \det[K(x_i,y_j)]_{i,j=1}^n
 =\frac1{n!}\int_{(0,\infty)^n}
 \det\left[\frac1{x_i+s_\ell^2}\right]_{i,\ell}
 \det\left[\frac1{y_j+s_\ell^2}\right]_{j,\ell}
 \prod_{\ell=1}^n\omega(s_\ell)\dd s_1\cdots\dd s_n.
\end{equation}
Absolute integrability follows by expanding into finitely many products
of the finite one-variable integrals. This is the usual
Andr\'eief--Heine identity used in positive determinant arguments
\cite[Section~2]{Zudilin}, with its elementary expansion proof here.
For $0<s_1<\cdots<s_n$, both Cauchy determinants are positive by the
Cauchy product formula~\eqref{eq:Cauchy}, applied to $x_i,s_j^2$
and $y_i,s_j^2$. Their product is unchanged by permuting the
$s_\ell$ and is strictly positive away from coincidences. This proves
the last assertion. It also proves that increasing the density increases
these determinants, even for different row and column nodes.
\end{proof}

\begin{remark}
The density in~\eqref{eq:K-integral} satisfies
$\omega(s)\sim\pi^2s^2/96$ at infinity; it has infinite total mass,
but the two resolvent factors make every kernel entry integrable.
Carreira's functional~\cite[Definition~24(B) and Lemma~25]{Carreira}
uses a positive density obtained from derivatives of $1/\sinh(\pi s)$,
with exponential decay and finite polynomial moments. The representation
used here is derived directly, and assumes no polynomial moments of
$\omega(s)\dd s$.
\end{remark}

\subsection{A uniform comparison with Cauchy determinants}

The positive representation makes nonvanishing transparent, but its
density is inconvenient for estimating a large determinant. We compare
that density with simpler ones whose determinants have explicit product
formulas. The comparison changes the logarithm by only $O(n)$, so it
preserves the $n^2$ scale relevant to the obstruction. Positivity of the
determinant integrand is what makes this comparison legitimate.

For positive nodes $a_i,b_j$, write
\begin{equation}\label{eq:Cauchy}
 C(a,b)_{ij}=\frac1{a_i+b_j},\qquad
 \det C(a,b)=
 \frac{\prod_{i<j}(a_j-a_i)(b_j-b_i)}{\prod_{i,j}(a_i+b_j)}.
\end{equation}
The product formula follows, for example, by multiplying the determinant
by its displayed denominator: the numerator is alternating separately in
the two node sets, hence contains their Vandermonde factors, and comparison
of degrees and one residue fixes the remaining constant to one.
The constants below have three roles: $M$ bounds the density factor
$q$, $\tau_0$ fixes a lower comparison that vanishes quadratically at
zero, and $c_0$ includes the factor $\pi/2$ from integrating the
comparison kernel. Set
\begin{equation}\label{eq:constants}
 c_0=\frac{\pi^3}{192},\qquad M=1+\frac{2}{3\sqrt6},\qquad
 \tau_0=\frac1{\sqrt2\pi}.
\end{equation}

\begin{theorem}[Uniform Cauchy comparison]\label{thm:comparison}
For any two strictly increasing lists of positive integers $a_i,b_j$,
\begin{equation}\label{eq:comparison}
 \left(\frac{c_0}{(1+\tau_0)^2}\right)^n\det C(a,b)
 \leq\det[K(a_i^2,b_j^2)]
 \leq (Mc_0)^n\det C(a,b).
\end{equation}
In particular,
\begin{equation}\label{eq:uniform-log}
 \log\det[K(a_i^2,b_j^2)]=\log\det C(a,b)+O(n),
\end{equation}
with an absolute constant independent of all node positions and gaps.
\end{theorem}

\begin{proof}
The maximum of $1+A-2A^3$ for $0\leq A\leq1$ is $M$. Also
$\cosh v\geq1+v^2/2$ and~\eqref{eq:q-positive} imply
\begin{equation}\label{eq:density-comparison}
 \frac{s^2}{s^2+\tau_0^2}\leq1-\sech(2\pi s)\leq q(s)\leq M.
\end{equation}
The elementary rational integrals needed for the comparison are
\begin{align}
 \int_0^\infty\frac{s^2\dd s}{(s^2+a^2)(s^2+b^2)}
   &=\frac\pi{2(a+b)},\label{eq:rational-integral1}\\
 \int_0^\infty\frac{s^4\dd s}
 {(s^2+\tau^2)(s^2+a^2)(s^2+b^2)}
   &=\frac\pi2 R_\tau(a,b),\label{eq:rational-integral2}\\
 R_\tau(a,b)&=\frac{ab+\tau(a+b)}{(a+b)(a+\tau)(b+\tau)}.
 \label{eq:Rtau}
\end{align}
To check the second formula, first divide
$s^4/((s^2+a^2)(s^2+b^2))$ into
$[a^2s^2/(s^2+a^2)-b^2s^2/(s^2+b^2)]/(a^2-b^2)$,
and then use the first formula with the extra factor $s^2+\tau^2$.
Coincident parameters follow by continuity. The first formula follows
by partial fractions and $\int_0^\infty(s^2+a^2)^{-1}\dd s=\pi/(2a)$.

Applying the positive determinant integral~\eqref{eq:andreief} to the
density bounds gives
\begin{equation}\label{eq:regularized-bounds}
 c_0^n\det[R_{\tau_0}(a_i,b_j)]
 \leq\det[K(a_i^2,b_j^2)]\leq(Mc_0)^n\det C(a,b).
\end{equation}
This uses the sign of the determinant integrand; an entrywise matrix
comparison alone would not suffice.

We now evaluate the regularized determinant exactly:
\begin{equation}\label{eq:Rtau-det}
 \begin{split}
 \det[R_\tau(a_i,b_j)]={}&\det C(a,b)
 \prod_i\frac{a_i}{a_i+\tau}\prod_j\frac{b_j}{b_j+\tau}\\
 &\times\left(1+\tau\sum_i\frac1{a_i}+\tau\sum_j\frac1{b_j}\right).
 \end{split}
\end{equation}
Let $u_i=1/a_i$, $v_j=1/b_j$. The entrywise factorization
\[
 R_\tau(a_i,b_j)=\frac{a_i}{a_i+\tau}
  \left(\frac1{a_i+b_j}+\frac{\tau}{a_i b_j}\right)
  \frac{b_j}{b_j+\tau}
\]
reduces the calculation to $\det(C+\tau uv^t)$.
The rank-one identity
$\det(C+\tau uv^t)=\det C(1+\tau v^tC^{-1}u)$ leaves one scalar
to compute. If $w=C^{-1}u$ and
$F(x)=\sum_jw_j/(x+b_j)$, then $F(a_i)=1/a_i$. Therefore
\[
 \frac1x-F(x)=\frac1x\,
 \frac{\prod_i(1-x/a_i)}{\prod_j(1+x/b_j)}.
\]
Indeed, after multiplication by $x\prod_j(x+b_j)$, the numerator has
degree at most $n$ and vanishes at every $a_i$; its factor is fixed by
the residue one at zero. Comparing constant terms at zero gives
$v^tC^{-1}u=F(0)=\sum_i1/a_i+\sum_j1/b_j$.
The rank-one determinant identity proves~\eqref{eq:Rtau-det}.
Finally each integer node is at least one, so the products are at least
$(1+\tau_0)^{-2n}$ and the last factor is at least one.
This proves~\eqref{eq:comparison} and~\eqref{eq:uniform-log}.
\end{proof}

\section{First construction: the arithmetic obstruction}\label{sec:principal}

We now determine part of the arithmetic cost in~\eqref{eq:budget}.
For each prime in a suitable interval, the fourth-power tails contribute
a diagonal residue that the cubic tails cannot cancel. A coefficient
of the determinant records this residue. One such coefficient per prime
is enough for a lower bound on the actual primitive multiplier.

\subsection{A coefficient with forced valuation}

For an odd prime $p$ put $\chi_4(p)=(-1)^{(p-1)/2}$.
Writing $M_0$ for the largest node, the interval $2M_0<p<4M_0$
below is chosen to isolate a single
denominator $p$ in each beta partial sum, while keeping all factors
from differences of nodes invertible modulo $p$. Multiplication by
$p^4$ then kills the cubic terms and detects only the fourth-power
diagonal terms. We call $a_i$ a high node when $4a_i>p$.

\begin{lemma}\label{lem:principal-local}
Let $a_1<\cdots<a_n$ be positive integers with maximum $M_0$, and set
$Q(X)=\det(XI_n+R(a))$. If $p\geq5$ and $2M_0<p<4M_0$, put
$r_p=\#\{i:4a_i>p\}$. Then
\begin{equation}\label{eq:principal-local}
 \vp\bigl(\coeff{n-r_p}{Q}\bigr)=-4r_p,
 \qquad \vp\bigl(\kappa(Q)\bigr)\geq4r_p,
\end{equation}
where $\PP(Q)=\kappa(Q)Q$.
\end{lemma}

\begin{proof}
The quantities $a_i$, $a_i+a_j$ and $a_i-a_j$ for $i\ne j$ are $p$-units.
Among odd denominators below $4M_0$, the only possible multiple of $p$
is $p$ itself, since $3p>4M_0$. It appears in $S_\ell(a_i)$ exactly when
$4a_i>p$, with sign $\chi_4(p)$. Consequently $p^4R(a)$ is $p$-integral
and its reduction is diagonal:
\begin{equation}\label{eq:principal-residue}
 p^4R(a)_{ij}\equiv
 \begin{cases}
 -\chi_4(p),&i=j,\ 4a_i>p,\\
 0,&\text{otherwise}
 \end{cases}\pmod p.
\end{equation}
Expanding $\det(XI+R)$ by the set of diagonal entries from which $X$
is selected shows that the coefficient of $X^{n-r_p}$ is the sum
of all principal minors of $R$ of size $r_p$.
After multiplication by $p^{4r_p}$, exactly one such
minor survives modulo $p$: that using all the high nodes. Its residue is
$(-\chi_4(p))^{r_p}\ne0$. For $r_p=0$ this is the leading coefficient
one. This proves the equality in~\eqref{eq:principal-local};
Lemma~\ref{lem:primitive} proves the inequality for the actual primitive
multiplier.
\end{proof}

\begin{remark}[A three-node example]\label{rem:prime-example}
Take nodes $1,2,3$ and $p=7$. The high nodes in the lemma are $2,3$,
and $\chi_4(7)=-1$, so
\[
 7^4 R\equiv\begin{pmatrix}0&0&0\\0&1&0\\0&0&1\end{pmatrix}\pmod7.
\]
The coefficient of $X$ in $\det(XI_3+R)$ is the sum of the three
principal $2\times2$ minors:
\[
 [X]Q=(R_{11}R_{22}-R_{12}^2)
       +(R_{11}R_{33}-R_{13}^2)
       +(R_{22}R_{33}-R_{23}^2).
\]
After multiplication by $7^8$, only the final minor survives modulo
$7$, with residue $1$. Thus $v_7([X]Q)=-8$, and the primitive multiplier
must contain at least $7^8$. This is the role of the uniquely surviving
minor: it rules out cancellation in the sum defining that coefficient.
Other coefficients may force additional powers; a lower bound is enough.
\end{remark}

Since $Q$ is monic, the primitive multiplier is integral. Thus the local
inequalities can be summed without any negative contribution from omitted
primes:
\begin{equation}\label{eq:principal-sum}
 \log\kappa(Q)\geq
 4\sum_{2M_0<p<4M_0}r_p\log p
\end{equation}
for $M_0$ sufficiently large. This is a lower bound after all coefficient
content has been removed.

\subsection{Proportional offsets}

We now compare two functions of the relative offset $m/n$:
$\mathcal A$ measures the forced denominator cost, and $\mathcal D$
measures the real decay. Their difference will be the lower bound
for the primitive value. This subsection treats finite limiting
offsets; the next handles $m/n\to\infty$ separately so that no
uniformity in an unbounded parameter is assumed.

Take $a_i=m+i$, $M_0=m+n$, and suppose first that $m/n\to\rho<\infty$.
Reverse the finite sum in~\eqref{eq:principal-sum}. Since the endpoints
$2M_0$ and $4a_i$ are even and exceed two, endpoint conventions for primes
have no effect. The prime number theorem, uniformly for
$2M_0\leq x\leq4M_0$, gives
\begin{align}
 4\sum_{2M_0<p<4M_0}r_p\log p
 &=4\sum_{a_i>M_0/2}
       \bigl(\vartheta(4a_i)-\vartheta(2M_0)\bigr)\notag\\
 &=16\sum_{a_i>M_0/2}(a_i-M_0/2)+o(n^2)\notag\\
 &=\mathcal A(\rho)n^2+o(n^2),\label{eq:A-rho}
\end{align}
where
\begin{equation}\label{eq:A-def}
 \mathcal A(\rho)=
 \begin{cases}2(1+\rho)^2,&0\leq\rho\leq1,\\8\rho,&\rho\geq1.
 \end{cases}
\end{equation}
Reversing the prime sum counts, for each node $a_i$, the primes from
$2M_0$ up to $4a_i$. With $a_i/n\to\rho+t$, the resulting
Riemann sum is $16\int_0^1(\rho+t-(\rho+1)/2)_+\dd t$, where
$u_+=\max(u,0)$. When $\rho\leq1$, only the upper part of the node
interval contributes; when $\rho\geq1$, every node contributes.
These are the two cases in~\eqref{eq:A-def}.
Uniformity of the prime-number-theorem error follows from
$\sup_{2M_0\leq x\leq4M_0}|\vartheta(x)-x|=o(M_0)$ and
$M_0=O(n)$.

For the real estimate, the differences $a_j-a_i=j-i$ do not depend
on the offset. The Vandermonde product is
$\prod_{i<j}(j-i)=\prod_{j=1}^{n-1}j!$, so the Cauchy product is
\begin{equation}\label{eq:C-shift}
 \det C(a,a)=\frac{\prod_{j=1}^{n-1}(j!)^2}
                    {\prod_{i,j=1}^n(2m+i+j)}.
\end{equation}
Stirling's formula, followed by a Riemann sum for the denominator, gives
\begin{equation}\label{eq:D-rho}
 \log\det C(a,a)=-\mathcal D(\rho)n^2+o(n^2),\qquad
 \mathcal D(\rho)=\frac32+\int_0^1\int_0^1\log(2\rho+x+y)\dd x\dd y.
\end{equation}
For clarity,
$2\sum_{j<n}\log(j!)=n^2\log n-\frac32n^2+O(n\log(n+1))$.
At $\rho=0$ the logarithmic corner singularity is integrable and its
Riemann sums converge. Indeed, counting at most $k$ grid points with
$i+j=k$ bounds the contribution from $i+j\leq\varepsilon n$ by
$O(\varepsilon^2(1+|\log\varepsilon|))+O(\log(n+1)/n)$.
A positive offset tending to zero is handled by
monotonicity between offsets zero and any fixed $\varepsilon>0$.
Integrating explicitly yields
\begin{equation}\label{eq:D-explicit}
 \begin{split}
 \mathcal D(\rho)=\frac12\bigl(& (2\rho+2)^2\log(2\rho+2)
 -2(2\rho+1)^2\log(2\rho+1)\\
 &+(2\rho)^2\log(2\rho)\bigr),
 \end{split}
\end{equation}
where $0^2\log0=0$.

It remains to show that shifting cannot make real decay overtake
the arithmetic cost. Differentiating the integral formula shows
that the derivative of $\mathcal D$ is decreasing, and
\begin{equation}\label{eq:D-derivative}
 \mathcal D'(\rho)=2\int_0^1\int_0^1
       \frac{\dd x\dd y}{2\rho+x+y}
 \leq\mathcal D'(0)=4\log2<4.
\end{equation}
On both pieces of~\eqref{eq:A-def}, $\mathcal A'(\rho)\geq4$.
It follows that
\begin{equation}\label{eq:A-minus-D}
 \mathcal A(\rho)-\mathcal D(\rho)
 \geq\mathcal A(0)-\mathcal D(0)=2-2\log2.
\end{equation}
Combining~\eqref{eq:principal-sum}, \eqref{eq:A-rho},
\eqref{eq:D-rho} and Theorem~\ref{thm:comparison} proves the stronger
proportional-offset statement
\begin{equation}\label{eq:proportional-result}
 \log Z_{n,m}(\theta)
 \geq\bigl(\mathcal A(\rho)-\mathcal D(\rho)\bigr)n^2+o(n^2).
\end{equation}

\subsection{Offsets larger than the matrix size}

Suppose now that $m/n\to\infty$. All nodes exceed $M_0/2$ eventually.
The relevant scale is now $mn$: the prime cost is of this order,
whereas the logarithm of the real determinant has size at most
$O(n^2\log(1+m/n))=o(mn)$.
The same reversal of sums gives
\begin{align}
 \log H_{n,m}
 &\geq16\sum_{i=1}^n\bigl(m+i-(m+n)/2\bigr)+o(nM_0)\notag\\
 &=8mn+8n+o(mn).\label{eq:large-m-arithmetic}
\end{align}
The relative prime-number-theorem error is still uniform on
$[2M_0,4M_0]$, so this argument imposes no upper bound on $m$.
By~\eqref{eq:C-shift} and the numerator estimate just used,
\begin{equation}\label{eq:large-m-real}
 \log\det C(a,a)
 \geq-n^2\bigl(\tfrac32+\log(2(1+m/n))\bigr)-O(n\log(n+1))
 =-o(mn).
\end{equation}
The Cauchy comparison costs only $O(n)$, whence
\begin{equation}\label{eq:large-m-result}
 \log Z_{n,m}(\theta)\geq(8-o(1))mn.
\end{equation}
Every sequence of ratios $m_n/n$ has subsequences converging in the
compactified interval $[0,\infty]$. If the liminf asserted in
Theorem~\ref{thm:principal} failed, a subsequence witnessing that failure
would have a further subsequence with such a limit. Equations
\eqref{eq:proportional-result}--\eqref{eq:large-m-result} exclude every
possible limit. Positivity follows from Proposition~\ref{prop:kernel}.
This completes the proof of Theorem~\ref{thm:principal}.

\section{Second construction: a moment determinant}\label{sec:gram}

The second design begins with a positive double integral. Its monomial
integrals involve fourth-power beta tails on the diagonal and cubic
tails off the diagonal. Expanding an even power of $u-v$ makes the
unwanted $\beta(3)$ terms cancel, leaving rational linear forms in
$\theta$. Moreover $y=(u-v)^2$ lies in $[0,1]$, so its powers form a
positive moment sequence on a short interval. This gives both a
nonvanishing Gram determinant and an effective estimate of its real size.

\subsection{The moments and their determinant}\label{sec:gram-construction}

For a polynomial $P(u,v)$, define the finite integral
\begin{equation}\label{eq:J}
 J(P)=\int_0^1\int_0^1
       \frac{P(u,v)\log^2(uv)}{\sqrt{uv}(1+uv)}\dd u\dd v,
 \qquad \mu_m=\frac1{96}J((u-v)^{2m}).
\end{equation}
The weight is integrable and strictly positive on the open square.
Since $y=(u-v)^2$ lies in $[0,1]$, the $\mu_m$ are moments of a
positive measure on that interval. The matrix $[\mu_{i+j}]$ is the
Gram matrix of $1,y,\ldots,y^{n-1}$; its determinant is positive
because no nonzero polynomial vanishes throughout the interval.
Positive-weight representations and Pad\'e approximants for Hurwitz zeta
values are developed by Pr\'evost~\cite{Prevost}. Rational modifications
of moment functionals are part of the classical orthogonal-polynomial
framework; Krattenthaler~\cite{Krattenthaler} gives determinant formulas
including such modifications and repeated parameters.

To identify the polynomial in $\theta$, expand $(u-v)^{2m}$.
The middle term contributes
$\binom{2m}{m}(\theta-T_4(m))$; paired off-diagonal terms have
rational values because their cubic beta tails cancel. The factor
$1/96$ normalizes $\mu_0=\theta$. The calculation below proves
$\mu_m=C_m\theta+B_m$, with the following rational coefficients:
\begin{align}
 T_\ell(j)&=\sum_{k=0}^{j-1}\frac{(-1)^k}{(2k+1)^\ell},
 &C_m&=\binom{2m}{m},\label{eq:T}\\
 B_m&=-C_mT_4(m)+\frac16
 \sum_{\substack{0\leq j\leq 2m\\j\ne m}}
       \binom{2m}{j}\frac{T_3(j)}{j-m},\label{eq:B}\\
 Q_n^{\Delta}(X)&=\det[C_{i+j}X+B_{i+j}]_{0\leq i,j<n},
 &Z_n^{\Delta}&=\kappa_n Q_n^{\Delta}=\PP(Q_n^{\Delta}),\quad \kappa_n>0.
 \label{eq:delta-Q}
\end{align}
Here $T_\ell(j)$ truncates the beta series after $j$ individual terms,
so $S_\ell(a)=T_\ell(2a)$. The central binomial coefficient $C_m$
comes from the middle term just described; $B_m$ collects the finite
tail corrections and the off-diagonal terms. Empty sums are zero.
In particular $B_0=0$ and $B_1=-149/81$.

\begin{lemma}\label{lem:moment-formula}
With $C_m,B_m$ defined in~\eqref{eq:T}--\eqref{eq:B},
\begin{equation}\label{eq:mu}
 \mu_m=C_m\theta+B_m.
\end{equation}
\end{lemma}

\begin{proof}
First compute the integral of one monomial; then expand $(u-v)^{2m}$.
The identity
\[
 \int_0^1\int_0^1u^{a+t-1}v^{b+t-1}\dd u\dd v
          =\frac1{(a+t)(b+t)}
\]
shows why two common-parameter derivatives insert $\log^2(uv)$.
Expand $(1+uv)^{-1}$ as a geometric series. For $a,b>0$, the
log-squared integral in each term is therefore the second derivative
of the right side at $t=0$. For $a\ne b$ this derivative is
\[
 \frac2{b-a}\left(\frac1{a^3}-\frac1{b^3}\right),
\]
and for $a=b$ it is $6/a^4$. The sum of the absolute integrals of the
terms is finite (their tails are $O(k^{-4})$), justifying termwise
integration. Substitution of $a=k+i+1/2$, $b=k+j+1/2$ gives
\begin{align}
 \frac{J(u^iv^i)}{96}&=(-1)^i(\theta-T_4(i)),\label{eq:mono-diag}\\
 \frac{J(u^iv^j)}{96}&=-\frac{(-1)^i}{6}
                 \frac{T_3(i)-T_3(j)}{j-i}\quad
                 (i\ne j,\ i+j\text{ even}).\label{eq:mono-off}
\end{align}
In~\eqref{eq:mono-off}, the two cubic tails have the same sign
$(-1)^i=(-1)^j$, so their $\beta(3)$ terms cancel.
Now expand
$\mu_m=\sum_{j=0}^{2m}(-1)^j\binom{2m}{j}J(u^jv^{2m-j})/96$.
Its diagonal term contributes
$\binom{2m}{m}(\theta-T_4(m))$ by~\eqref{eq:mono-diag}.
For $j\ne m$, the sign $(-1)^j$ cancels the corresponding sign
in~\eqref{eq:mono-off}. The off-diagonal sum becomes
\[
 -\frac1{12}\sum_{\substack{0\leq j\leq2m\\j\ne m}}\binom{2m}{j}
       \frac{T_3(j)-T_3(2m-j)}{m-j}.
\]
Changing $j$ to $2m-j$ in the second part makes the two parts equal.
The result is the sum in~\eqref{eq:B}, proving~\eqref{eq:mu}.
\end{proof}

All coefficient statements use the explicitly defined polynomial
$C_mX+B_m$. No independence or irrationality of $\theta$ is assumed to
identify its coefficients.

\subsection{Positive measure and full degree}

On the half-square $u>v$, set $u=v+s$. The ranges are
$0<s<1$ and $0<v<1-s$. Symmetry doubles this integral, and then
$y=s^2$ gives $2\dd s=\dd y/s$. Thus $\mu_m$ is the $m$th
moment of a measure on $[0,1]$ with density
\begin{equation}\label{eq:rho}
 \rho(y)=\frac1{96s}\int_0^{1-s}
 \frac{\log^2(v(v+s))}{\sqrt{v(v+s)}(1+v(v+s))}\dd v,
 \qquad 0<y<1.
\end{equation}
This density is positive throughout $(0,1)$ and has total mass
$\mu_0=\theta$. Consequently
\begin{equation}\label{eq:Dpositive}
 D_n:=Q_n^{\Delta}(\theta)=\det[\mu_{i+j}]_{0\leq i,j<n}>0
\end{equation}
for every $n$: a nonzero polynomial has positive squared norm.

\begin{lemma}\label{lem:binomial-det}
For $n\geq1$,
\begin{equation}\label{eq:binomial-det}
 \det[C_{i+j}]_{0\leq i,j<n}=2^{n-1}.
\end{equation}
Thus $Q_n^{\Delta}$ has exact degree $n$ and leading coefficient $2^{n-1}$.
\end{lemma}

\begin{proof}
The $C_m$ are the moments of the probability arcsine measure
$\dd y/(\pi\sqrt{y(4-y)})$ on $[0,4]$. This follows either by the beta
integral or by the substitution $y=2+2\cos t$, $0<t<\pi$, and the
binomial expansion of $(2+2\cos t)^m$.
Let $\mathsf T_j$ denote the Chebyshev polynomial satisfying
$\mathsf T_j(\cos t)=\cos(jt)$. The monic integer polynomials
\[
 c_0(y)=1,\qquad c_j(y)=2\mathsf T_j((y-2)/2)\quad(j\geq1)
\]
are orthogonal for this measure, with squared norms $1,2,2,\ldots$.
Monicity and integrality follow from $c_1=y-2$, $c_2=(y-2)^2-2$ and
the Chebyshev recurrence for $j\geq2$. The coefficient change from
monomials has determinant one. The moment determinant is therefore the
product of these norms, giving~\eqref{eq:binomial-det}.
\end{proof}

\subsection{Matching upper and lower real exponents}

The upper estimate needs monic polynomials uniformly small on $[0,1]$;
Chebyshev polynomials supply them. For the lower estimate we compare
the measure with $(1-y)\dd y$, whose monic orthogonal norms can be
computed exactly. Changing from monomials to any monic polynomial
basis uses a triangular matrix with determinant one, so it preserves
the Gram determinant in both calculations.

\begin{proposition}\label{prop:gram-real}
For all $n\geq1$,
\begin{equation}\label{eq:gram-bounds}
 \frac1{6144^n n!\,2^{2n^2+n}}
 \leq D_n\leq \theta^n4^{-(n-1)^2}.
\end{equation}
In particular,
\begin{equation}\label{eq:gram-rate}
 \log D_n=-(\log4)n^2+O(n\log(n+1)).
\end{equation}
\end{proposition}

\begin{proof}
For $j\geq1$ the monic shifted Chebyshev polynomial on $[0,1]$ has
supremum norm $2/4^j$. Change the monomial basis of the Gram matrix
to these monic polynomials and apply Hadamard's inequality. Since the
measure has total mass $\theta$,
\[
 D_n\leq\theta^n\prod_{j=1}^{n-1}(2/4^j)^2
       =\theta^n4^{-(n-1)^2}.
\]

For the lower bound restrict the integral in~\eqref{eq:rho} to
$v\in[(1-s)/4,(1-s)/2]$. There $v(v+s)\leq1/2$, so
$-\log(v(v+s))\geq\log2\geq1/2$, and the denominator of the
integrand is at most 2. The integrand is at least $1/8$, giving
\begin{equation}\label{eq:rho-lower}
 \rho(y)\geq\frac{1-s}{3072s}\geq\frac{1-y}{6144}
 \quad(0<s<1).
\end{equation}
The second inequality is $2\geq s+s^2$. Ordering the positive definite
Gram matrices therefore gives
\begin{equation}\label{eq:Jacobi-det}
 D_n\geq6144^{-n}
 \det\left[\frac1{(i+j+1)(i+j+2)}\right]_{0\leq i,j<n}
 =6144^{-n}\prod_{j=0}^{n-1}
    \frac1{(2j+2)\binom{2j+1}{j}^2}.
\end{equation}
Here determinant monotonicity follows by conjugating the matrix
inequality by the inverse square root of the smaller positive matrix.

We include the norm calculation in the last equality. For weight $1-y$
on $[0,1]$, the monic orthogonal polynomial of degree $j$ is
\[
 p_j(y)=\frac{(-1)^j(j+1)!}{(2j+1)!(1-y)}
         \frac{\mathrm d^j}{\mathrm dy^j}
             \bigl(y^j(1-y)^{j+1}\bigr).
\]
The numerator is divisible by $1-y$. Integration by parts $j$ times
shows orthogonality to degrees below $j$; every boundary term vanishes.
The same calculation with $p_j$ gives its squared norm
\[
 \frac{j!(j+1)!}{(2j+1)!}
       \int_0^1 y^j(1-y)^{j+1}\dd y
 =\frac{j!^2(j+1)!^2}{(2j+1)!(2j+2)!}
 =\frac1{(2j+2)\binom{2j+1}{j}^2}.
\]
The Gram determinant is the product of these monic norms.
Finally $\binom{2j+1}{j}\leq2^{2j+1}$ in~\eqref{eq:Jacobi-det}
gives the lower bound in~\eqref{eq:gram-bounds}. Taking logarithms,
using $0<\theta<1$ and Stirling's formula, proves~\eqref{eq:gram-rate}.
\end{proof}

\section{Second construction: the arithmetic obstruction}\label{sec:delta-arithmetic}

Here the prime-denominator contributions first appear at a specific
moment index. Because a moment matrix has entry index $i+j$, their
residues occupy a corner block with a nonzero anti-diagonal. Its
determinant supplies the noncancelling coefficient, just as the single
surviving minor did in Remark~\ref{rem:prime-example}. In this family
we also bound every coefficient's denominator at the same prime, which
turns a lower bound into an exact formula.

Fix $n\geq1$ and a prime $p>4n$. All moment indices used in
$Q_n^{\Delta}$ satisfy $0\leq m\leq2n-2$. In particular $2m<p$,
$|j-m|<p$ for the indices in~\eqref{eq:B}, and $T_4(m)$ is $p$-integral.
In $T_3(j)$ the only possible odd denominator divisible by $p$ is $p$
itself: the largest odd denominator is at most $8n-9<2p$, whereas the
next odd multiple is $3p$. It appears when $j\geq(p+1)/2$.

\subsection{The first residue and its block}

It follows that $p^3B_m$ is $p$-integral and
\begin{equation}\label{eq:Bm-residue}
 p^3B_m\equiv\frac{\chi_4(p)}6
  \sum_{j=h}^{2m}\binom{2m}{j}\frac1{j-m}\pmod p,
 \qquad h=\frac{p+1}{2}.
\end{equation}
An empty sum is zero; in the range of this sum $j>m$, so no zero
denominator occurs. Put
\begin{equation}\label{eq:m0-r}
 m_0=\left\lfloor\frac p4\right\rfloor+1,\qquad
 r=\max\{0,2n-1-m_0\}.
\end{equation}
Here $m_0$ is the first moment index for which the sum in
\eqref{eq:Bm-residue} can be nonempty: $2m\geq(p+1)/2$ is
equivalent to $m\geq\lfloor p/4\rfloor+1$.
For $m<m_0$ no cubic tail contains $p$, so $B_m$ is actually $p$-integral.
If $m_0\leq2n-2$, the first nonzero residue in~\eqref{eq:Bm-residue},
apart from its factor $\chi_4(p)/6$, is
\begin{equation}\label{eq:first-residue}
 \begin{cases}
 1/m_0,&p=4m_0-1,\\[2pt]
 (2m_0-1)(m_0+1)/(m_0(m_0-1)),&p=4m_0-3.
 \end{cases}
\end{equation}
In the first case the sum has only $j=2m_0$; in the second it has
$j=2m_0-1,2m_0$, giving $2m_0/(m_0-1)+1/m_0$.
These residues are nonzero modulo $p$. Indeed, in the nonempty case
$m_0\geq3$, and all the displayed positive numerator and denominator
factors are less than $p$.

If $r=0$, every residue vanishes. Otherwise, row $i$ can contain
a nonzero residue only if $i+j\geq m_0$ for some $j\leq n-1$,
hence $i\geq m_0-(n-1)=n-r$; the same applies to columns.
The residue matrix $[p^3B_{i+j}\bmod p]_{i,j<n}$ is therefore
supported in the bottom-right $r\times r$ block.
That block is anti-triangular with the
nonzero first residue on its anti-diagonal, so it has rank $r$.
More strongly, every entry outside it is $p$-integral before multiplication
by $p^3$: its moment index is below $m_0$.

\subsection{The exact primitive exponent}

\begin{proposition}[Exact large-prime cost]\label{prop:gram-prime-cost}
For every $n\geq1$ and every prime $p>4n$, the primitive multiplier
$\kappa_n$ of $Q_n^\Delta$ satisfies
\begin{equation}\label{eq:delta-valuation}
 \vp(\kappa_n)=3\max\left\{0,\,2n-2-\left\lfloor\frac p4\right\rfloor\right\}.
\end{equation}
\end{proposition}

\begin{proof}[Proof of \eqref{eq:delta-valuation}]
At most $r$ entries of a permutation term can have negative valuation,
because only the bottom $r$ rows have such entries. Each entry has
coefficient valuations at least $-3$, so every coefficient of
$Q_n^{\Delta}$ has valuation at least $-3r$.

To see that the bound is attained, suppose $r>0$, let $L$ be the
$r\times r$ residue block just described, and put $s=n-r$.
Multiply the bottom $r$ rows of the pencil by $p^3$. The new matrix
has $p$-integral polynomial entries, and its reduction has the form
\[
 \begin{pmatrix}
 [C_{i+j}X+B_{i+j}]_{0\leq i,j<s}& *\\
 0&L
 \end{pmatrix}.
\]
The lower-left block vanishes because those entries were already
$p$-integral. In the lower-right block the integral $X$-coefficients
also vanish after multiplication by $p^3$. Taking this block-triangular
determinant proves the polynomial congruence
\begin{equation}\label{eq:block-factor}
 p^{3r}Q_n^{\Delta}(X)
 \equiv \det(L)\det[C_{i+j}X+B_{i+j}]_{0\leq i,j<s}\pmod p.
\end{equation}
The right-hand coefficient of $X^s$ is $\det(L)2^{s-1}$ by
Lemma~\ref{lem:binomial-det}, and is a $p$-unit. Here $s\geq1$,
since $p>4n$. Thus the minimum coefficient valuation is exactly $-3r$.
If $r=0$, all entries are $p$-integral and the leading coefficient
$2^{n-1}$ is already a $p$-unit. Lemma~\ref{lem:primitive} now proves
\eqref{eq:delta-valuation} in all cases.
\end{proof}

This argument controls the actual primitive multiplier. Any coefficient
gcd is already included in the minimum-valuation calculation; it cannot
remove the forced exponent.

\subsection{Summing the prime cost}

By Lemma~\ref{lem:primitive} and the leading coefficient $2^{n-1}$,
$\vp(\kappa_n)\geq0$ at every odd prime, while
$v_2(\kappa_n)\geq-(n-1)$. Hence
\begin{equation}\label{eq:kappa-sum}
 \log\kappa_n\geq
 3\sum_{4n<p<8n-8}\left(2n-2-\left\lfloor\frac p4\right\rfloor\right)
                   \log p-(n-1)\log2.
\end{equation}
Changing the endpoints by a bounded amount and removing floors costs
$O(n\log(n+1))$, which is negligible on the quadratic scale. Partial
summation and~\eqref{eq:PNT} give
\begin{align}
 3\sum_{4n<p<8n-8}\left(2n-2-\left\lfloor\frac p4\right\rfloor\right)
                   \log p
 &=3\int_{4n}^{8n}(2n-t/4)\dd\vartheta(t)+O(n\log(n+1))\notag\\
 &=3\int_{4n}^{8n}(2n-t/4)\dd t+o(n^2)
  =6n^2+o(n^2).\label{eq:six}
\end{align}
For the replacement of $\dd\vartheta$ by $\dd t$, integration by parts
bounds the error using $\sup_{4n\leq t\leq8n}|\vartheta(t)-t|=o(n)$,
the $O(n)$ endpoint weight and its bounded derivative.
Thus $\log\kappa_n\geq6n^2+o(n^2)$. Combining with
Proposition~\ref{prop:gram-real} proves~\eqref{eq:delta-result} and
completes the proof of Theorem~\ref{thm:delta}.

\section{Scope, evidence and the wider investigation}\label{sec:scope}

The two proofs establish actual primitive-value obstructions. The first
uses a monic pencil and a unique surviving principal minor; the second
uses a power-of-two leading coefficient and a full anti-triangular
residue block. In both cases the leading coefficient controls possible
negative valuations of the primitive multiplier at the primes not used
in the lower-bound sum. Without that control, a large denominator cost
alone would not prove the same conclusion.

These results concern the full pencils defined in
\eqref{eq:principal-Q} and~\eqref{eq:delta-Q}. Invertible rational basis
changes preserve their primitive polynomials, but different subspaces,
sparse designs, nonprincipal minors, polynomial weights, factors and
combinations of determinants are outside the two obstruction theorems.
The uniform real comparison in Theorem~\ref{thm:comparison} does apply
to nonprincipal designs; an arithmetic obstruction for such a design
would require additional information.

The accompanying research catalog separates proved limitations of defined
families from failed estimates, exact reductions to existing constructions,
finite experiments and incomplete constructions. In particular, some
nonprincipal designs have unfavorable proved denominator bounds without
an all-index bound on numerator content. An ordinary-moment first-variation
construction has a positive error and a proved real exponent, but no
all-index primitive arithmetic conclusion. Neither status is upgraded
to an impossibility statement by the present theorems.
The accompanying \href{run:beta4-investigations.pdf}{technical supplement}
develops selected further proofs and guides the reader through the
remaining constructions, with links to the catalog and source snapshots.

\subsection*{Computational checks and authorship}

This note was prepared with Codex. The archived record includes model
mathematical cross-checks, exact rational calculations and Arb interval
calculations. A reproducibility record accompanies this draft and states
which checks were rerun and how the archived outputs were preserved.
Such checks are finite evidence for identities and arithmetic formulas;
the asymptotic results are proved in the text. No proof-assistant
verification or human mathematical audit is claimed. The comparison with
prior work is recorded by source and version; no exhaustive priority claim
is made.

\appendix
\section{Including a common zero node}\label{app:zero}

The unregularized Cauchy matrix is unsuitable when both node lists contain
zero, but the regularized matrix in~\eqref{eq:Rtau} remains finite. For
$a_0=b_0=0$ and increasing positive remaining nodes, write $a_+,b_+$ for
the lists with zero removed. Continuity gives $R_\tau(0,0)=1/\tau$ and
\begin{equation}\label{eq:zero-Rdet}
 \det[R_\tau(a_i,b_j)]_{0\leq i,j<n}
 =\frac1\tau\prod_{i>0}\frac{a_i}{a_i+\tau}
              \prod_{j>0}\frac{b_j}{b_j+\tau}\det C(a_+,b_+).
\end{equation}
To prove it, multiply row $i$ by $a_i+\tau$ and column $j$ by $b_j+\tau$.
The resulting entry is $\tau+ab/(a+b)$, with continuous value $\tau$
in the zero row and column. Subtract the zero row from every other row
and expand at the zero column. Undoing the scalings gives
\eqref{eq:zero-Rdet}; an empty determinant is one.

The lower density bound~\eqref{eq:density-comparison} still applies.
For an upper bound finite at zero, note that
$1-\sech v\leq v^2/2$ for $v\geq0$: integrate
$\sech v\tanh v\leq v$. Therefore
\[
 q(s)\leq5(1-\sech(2\pi s))\leq10\pi^2s^2,
 \qquad q(s)\leq M.
\]
Combining the bounds gives
$q(s)\leq(M+10\pi^2\tau^2)s^2/(s^2+\tau^2)$ for every $\tau>0$.
Choose $\tau=\tau_0$; then
\begin{equation}\label{eq:zero-comparison}
 c_0^n\det R_{\tau_0}
 \leq\det[K(a_i^2,b_j^2)]
 \leq[c_0(M+5)]^n\det R_{\tau_0}.
\end{equation}
For positive integer remaining nodes,~\eqref{eq:zero-Rdet} gives
\begin{equation}\label{eq:zero-log}
 \log\det[K(a_i^2,b_j^2)]
 =\log\det C(a_+,b_+)+O(n),
\end{equation}
uniformly in the remaining node lists.

\begin{corollary}\label{cor:zero}
For the principal nodes $0,1,\ldots,n-1$, replace $\theta$ by $X$ in
\eqref{eq:K-off}--\eqref{eq:K-zero}, take the determinant, and let
$Z_n^0$ be its primitive representative. Then
\[
 \liminf_{n\to\infty}\frac{\log Z_n^0(\theta)}{n^2}
 \geq2-2\log2.
\]
\end{corollary}

\begin{proof}
The coefficient matrix of $X$ is $\diag(2,1,\ldots,1)$, so the leading
coefficient of the determinant is 2. For primes $2(n-1)<p<4(n-1)$,
the unique high-node minor from Lemma~\ref{lem:principal-local} is
unchanged; its complementary diagonal coefficient contributes a factor
2, a $p$-unit. Thus the same coefficient valuation $-4r_p$ is attained.
All odd prime valuations of the primitive multiplier are nonnegative,
and the only possible negative contribution is at least $-\log2$.
The prime sum is $2(n-1)^2+o(n^2)$. By~\eqref{eq:zero-log}, its positive
real determinant has logarithm $-2\log2\,(n-1)^2+o(n^2)$. Adding proves
the assertion.
\end{proof}

\begin{thebibliography}{9}

\bibitem{Brown}
F.~Brown,
\emph{Mellin transforms, transfinite diameter and rational approximations
of integrals}, preprint,
\href{https://arxiv.org/abs/2604.20741v1}{arXiv:2604.20741v1} (2026).

\bibitem{Carreira}
M.~C.~S.~Carreira,
\emph{Determinant constructions for the Dirichlet beta value at four},
preprint (2026),
\href{https://github.com/MarcosCarreira/catalan-designs/blob/974a6e6c7b4bce67538aa1b9bf5883b931789d1f/paper/beta4_determinants.pdf}{pinned PDF in the catalan-designs repository},
commit \texttt{974a6e6c7b4b} (full identifier in the accompanying source ledger).

\bibitem{Krattenthaler}
C.~Krattenthaler,
\emph{A determinant identity for moments of orthogonal polynomials that
implies Uvarov's formula for the orthogonal polynomials of rationally
related densities},
\href{https://arxiv.org/abs/2103.03969v1}{arXiv:2103.03969v1} (2021).

\bibitem{OpenAI}
OpenAI,
\emph{Catalan's constant is irrational}, manuscript dated
24 September 2026,
\href{https://github.com/openai/math/blob/adc7f1241b42e322a6451854ab7e4b4c146bf78a/preprints/Catalans-constant-is-irrational-September-24-2026/paper.pdf}{pinned PDF in the OpenAI math repository},
commit \texttt{adc7f1241b42} (full identifier in the source ledger).

\bibitem{Prevost}
M.~Pr\'evost,
\emph{Remainder Pad\'e approximants for the Hurwitz zeta function},
\href{https://arxiv.org/abs/1709.05389v1}{arXiv:1709.05389v1} (2017).

\bibitem{Zudilin}
W.~Zudilin,
\emph{A determinantal approach to irrationality},
\href{https://arxiv.org/abs/1507.05697v2}{arXiv:1507.05697v2} (2015).

\end{thebibliography}
\end{document}
